% SCA101 Lab 2_1A - What You Should See on a Real Husky
% Generated by guide_tools/latex/latexkit.py from the story script; this file is plain LaTeX - edit it directly.
% Build: pdflatex <this file>   (run it twice so the contents page fills in).  Needs scaguide.sty and figures/ next to it.
\documentclass[11pt,letterpaper]{article}
\usepackage{scaguide}
\guidesetup{
  headerleft={ChipWhisperer SCA101  -  Lab 2\_1A on a real Husky},
  headerright={What you should see},
  footerleft={Companion to: Lab 2\_1A (MAIN / HARDWARE) and the Lab 2\_1A Explorer},
  title={SCA101 Lab 2\_1A - What You Should See on a Real Husky},
  subject={Expected results for Lab 2\_1A and its Explorer notebook when VERSION is switched from SIMULATED to HARDWARE}}

\begin{document}

\guidecover{SCA101  -  LAB 2\_1A  (HARDWARE)}{What You Should See\\ on a Real Husky}{Expected results when you switch Lab 2\_1A and its Explorer from SIMULATED to HARDWARE}

Lab 2\_1A (\textbf{instruction power differences}) and its Explorer notebook run on \textbf{recorded traces} by
default. Setting \code{VERSION = 'HARDWARE'} in the Setup cell runs the same code on your own
\textbf{ChipWhisperer-Husky + CW313 + SAM4S} - and the traces look \textbf{different} from the recorded ones: longer,
a different loop speed, a different starting point. None of that is a fault. This guide shows what a healthy kit
produces, so you can tell ``it is working'' from ``something is wrong'', and how to measure a loop on a live trace.

\vspace{10pt}

\begin{guidecallout}{tip}{HOW TO USE THIS GUIDE}
All numbers here were \textbf{measured on one real kit} (Husky firmware 1.5.0, ChipWhisperer 6.0.0, SAM4S target,
Windows 11, Python 3.12, 25 dB gain). Yours will be close, not identical: look for the same \textbf{pattern}. The
captures behind every figure are bundled in
\code{jupyter/\bk{}courses/\bk{}sca101/\bk{}traces/\bk{}hardware\_\bk{}reference/\bk{}} so you can compare your own
against them.
\end{guidecallout}

\begin{guidecallout}{warn}{WHICH STEP 7?}
\textbf{Step 7 of the Explorer changed.} Version 1.0 of \code{Lab 2\_\bk{}1A Explorer - Student Guide.\bk{}pdf}
described an earlier Step 7 that counted pulses on the live trace with \code{find\_pulses()}; that does not work
reliably on real traces. \textbf{Version 1.1 of that guide and this guide match the notebook} (check the version on
the guide's cover). Steps 1 to 6 are unchanged.
\end{guidecallout}

\guidehthree{\textbf{What is in this guide}}

{\small\arrayrulecolor{guidegrid}
\rowcolors{2}{white}{guidezebra}
\begin{longtable}{|R{0.1159}|R{0.4638}|R{0.4203}|}
\hline
\rowcolor{guidenavy}\hdrcell{Section} & \hdrcell{What} & \hdrcell{Use it when} \\ \hline
\endhead
1 & What changes when VERSION = `HARDWARE' & You switch for the first time. \\ \hline
2 & The Setup cell: what it should print & You want to know it connected and configured correctly. \\ \hline
3 & The four traces: recorded next to a real chip & You open the first plots of the lab. \\ \hline
4 & Explorer Step 7: measure a loop you wrote yourself & You run Step 7 or want to know what numbers to expect. \\ \hline
5 & When it does not look like this & A plot is flat, saturated, or the capture times out. \\ \hline
6 & The bundled reference traces & You want to compare without hardware. \\ \hline
\end{longtable}}

\clearpage

\tableofcontents

\clearpage

\guidehone{1  What changes when VERSION = `HARDWARE'}

Open the notebook's first code cell and change one word, then choose \textbf{Kernel \textgreater{} Restart \& Run
All}:

\begin{guidecode}
VERSION = 'HARDWARE'
\end{guidecode}

Before you run it: the Husky kit assembled as in Lab 0 (CW313 with the \textbf{SAM4S} board, jumper \code{JP2} on
T-GPIO4, the \textbf{Pos} SMA cable from the Husky to the CW313), the Husky plugged in, and \textbf{no other notebook
holding it} - only one program can use the Husky at a time. The Setup cell then \textbf{builds and flashes} the lab
firmware (\code{simpleserial-base-lab2}) onto the target by itself.

\guidehthree{What is different in the data:}

{\small\arrayrulecolor{guidegrid}
\rowcolors{2}{white}{guidezebra}
\begin{longtable}{|R{0.3333}|R{0.3043}|R{0.3623}|}
\hline
\rowcolor{guidenavy}\hdrcell{} & \hdrcell{Recorded traces (SIMULATED)} & \hdrcell{Live traces (HARDWARE)} \\ \hline
\endhead
\textbf{Length of a trace} & 2000 samples & \textbf{5000 samples} (\codesmall{scope.adc.samples}) \\ \hline
\textbf{Samples per clock cycle} & 4 & \textbf{4} (target clock 7.36 MHz, ADC 29.45 MHz) - the same \\ \hline
\textbf{Amplitude} & about -0.3 to +0.2 & about \textbf{-0.4 to +0.34} (standard deviation about 0.16) \\ \hline
\textbf{Where the interesting code starts} & from about sample 0 & near \textbf{sample 410}: start-up code runs first \\ \hline
\textbf{Speed of the multiply loop (\codesmall{A *= 2;})} & \textbf{15} cycles (60 samples) per trip & \textbf{20} cycles (80 samples) per trip \\ \hline
\textbf{Speed of the divide loop (\codesmall{A /= 3;}, A = 0x7FFFFFFF)} & \textbf{22} cycles falling to \textbf{16} & \textbf{30} cycles falling to \textbf{23} \\ \hline
\textbf{How you change the code} & you cannot - fixed recordings & \textbf{you can}: Explorer Step 7 edits the C code, rebuilds, reflashes and captures again \\ \hline
\end{longtable}}

\vspace{4pt}

\begin{guidecallout}{tip}{GOOD TO KNOW}
\textbf{Why are the loop speeds different?} The recorded traces were taken from a different firmware build, so the
compiler laid the same C loop out differently. The loop still repeats a pattern, still scales with the count, and the
divide still gets faster as \code{A} shrinks - those ideas are what the lab teaches, and they all hold on the real
chip.
\end{guidecallout}

\Needspace{259pt}

\guidehone{2  The Setup cell: what it should print}

A healthy run prints the lines below (abridged here - the real output also shows the full firmware build). The numbers
that matter are in \textbf{bold} in the explanation underneath.

\begin{guidecode}
firmware working copy: ...\firmware\mcu\simpleserial-base-lab2
INFO: Found ChipWhisperer
scope.gain.gain          changed from 0       to 22
scope.gain.db            changed from 15.0    to 25.09
scope.adc.samples        changed from 131124  to 5000
scope.clock.clkgen_freq  changed from 0       to 7363636.36
scope.clock.adc_freq     changed from 0       to 29454545.45
scope.io.hs2             changed from None    to clkgen
scope.trace.capture.trigger_source  changed from trace trigger, rule #0  to firmware trigger
Hardware ready - baseline simpleserial-base flashed.
\end{guidecode}

{\small\arrayrulecolor{guidegrid}
\rowcolors{2}{white}{guidezebra}
\begin{longtable}{|R{0.3913}|R{0.6087}|}
\hline
\rowcolor{guidenavy}\hdrcell{Line} & \hdrcell{What it means} \\ \hline
\endhead
\codesmall{scope.\bk{}gain.\bk{}db .\bk{}.\bk{}.\bk{} to 25.\bk{}09} & The measurement amplifier is set to \textbf{25 dB}. This is the default and a good value; if your traces hit the +/-0.5 limit see section 5. \\ \hline
\codesmall{scope.\bk{}adc.\bk{}samples .\bk{}.\bk{}.\bk{} to 5000} & Every capture is \textbf{5000 samples} long. \\ \hline
\codesmall{scope.\bk{}clock.\bk{}clkgen\_\bk{}freq .\bk{}.\bk{}.\bk{} 7363636}, \codesmall{adc\_\bk{}freq .\bk{}.\bk{}.\bk{} 29454545} & The Husky generates the target's clock (\textbf{7.36 MHz}) and samples at \textbf{4 times that} (29.45 MHz): \textbf{4 samples per clock cycle}. \\ \hline
\codesmall{trigger\_\bk{}source .\bk{}.\bk{}.\bk{} firmware trigger} & The scope starts recording when the target firmware raises its trigger pin - which is why the board must hold the lab firmware. \\ \hline
\codesmall{Hardware ready - baseline simpleserial-base flashed.\bk{}} & The target was programmed. The next cell then prints \codesmall{trace length: 5000} and \codesmall{OK to continue!}. \\ \hline
\end{longtable}}

\vspace{4pt}

\begin{guidecallout}{ifnot}{IF IT DOESN'T}
\code{OSError: Unable to communicate with found ChipWhisperer} - another program holds the Husky (often a second
notebook). \textbf{Kernel \textgreater{} Shut Down All Kernels}, then run again. \code{Could not find ChipWhisperer} -
the cable or a USB rule (Linux guide). \code{Verify FAILED} while flashing - the wrong board is on the CW313, or it is
not seated; unplug and replug the Husky.
\end{guidecallout}

\Needspace{446pt}

\guidehone{3  The four traces: recorded next to a real chip}

The lab captures four short programs between two triggers and compares their power traces: \textbf{0} nothing,
\textbf{1} twenty multiplies in a row, \textbf{2} a loop of 20 multiplies, \textbf{3} a loop of 20 divides. The
picture shows the recorded trace (left) and a real capture (right) for each.

\guidefigure{figures/sections.png}{1.000}{Left: the recorded traces from Steps 1 - 6. Right: real captures of the same four programs. The live traces are denser (more samples) and start with about 400 samples of start-up activity before your code begins.}

{\small\arrayrulecolor{guidegrid}
\rowcolors{2}{white}{guidezebra}
\begin{longtable}{|R{0.3623}|R{0.6377}|}
\hline
\rowcolor{guidenavy}\hdrcell{On a real chip} & \hdrcell{What you should see} \\ \hline
\endhead
\textbf{Standard deviation of each trace} & about \textbf{0.16} (0.158 to 0.164 across the four) \\ \hline
\textbf{Largest value in a trace} & between about \textbf{-0.41} and \textbf{+0.34}. Anything stuck at +/-0.5 is clipping - section 5. \\ \hline
\textbf{Repeatability} & the same program captured twice differs by only about \textbf{0.005} per sample (mean absolute difference), so differences between programs are real, not noise. \\ \hline
\textbf{Sections 0 and 1} & no loop. They differ by only about \textbf{0.02} per sample before sample 400, and by about \textbf{0.11 to 0.12} per sample once the multiplies run. \\ \hline
\textbf{Section 2 (loop of 20 multiplies)} & a repeating pattern of \textbf{20 cycles = 80 samples}, starting at about sample \textbf{406} and lasting about \textbf{1630 samples} (20.4 trips). \\ \hline
\textbf{Section 3 (loop of 20 divides)} & a pattern whose length \textbf{keeps shrinking}, from about \textbf{30 cycles} to about \textbf{23} over the 20 trips - the divide gets faster as \codesmall{A} gets smaller. \\ \hline
\textbf{After your code} & another repeating pattern of \textbf{18 cycles = 72 samples} (it begins near sample 990 in section 0, 1490 in section 1 and 3110 in section 3). It is \emph{not} your loop. Step 7's measurement only accepts a loop that starts between sample 350 and 700, so it is not fooled by it. \\ \hline
\end{longtable}}

\vspace{4pt}

\guidehone{4  Explorer Step 7: measure a loop you wrote yourself}

Step 7 of the \textbf{Lab 2\_1A Explorer} is hardware only. You choose the C statement inside the loop and how many
times it runs; the notebook writes it into the firmware, rebuilds, reflashes (about 10 seconds), captures a trace and
measures the loop:

\begin{guidecode}
MY_OPERATION = "A *= 2;"      # the statement inside the loop
MY_COUNT     = 20             # how many trips
MY_A_START   = "0x2BAA"       # starting value of A
\end{guidecode}

\textbf{Why not use the pulse finder from Step 4?} It was tuned on the smooth recorded traces and finds only a couple
of pulses on a live trace. Step 7 uses a different measurement: it compares each block of \code{P} samples with the
next block. Inside a loop the blocks repeat, so the match is above 0.9 \textbf{at the loop's trip length and nowhere
else}.

\guidefigure{figures/loop_found.png}{0.899}{Step 7 with \code{A *= 2;} x 20 on a real chip. The green band is the loop the measurement found (sample 410 to 2036); the lower plot shows the trip length in each 300-sample window: a steady \textbf{20 cycles} while the loop runs (red), scattered values outside it (grey).}

\begin{guidecallout}{see}{YOU SHOULD SEE}
\code{every trip takes 80 samples = 20 clock cycles}, then \code{the loop runs from sample 410 to 2036 (1626 samples)
= about 20.\bk{}3 trips (you asked for 20)}. The largest swing printed above it is about 0.4; the ADC limit is 0.5.
\end{guidecallout}

\guidehthree{\textbf{What we measured for each experiment} (one real kit; your numbers should match to within a cycle or so):}

{\small\arrayrulecolor{guidegrid}
\rowcolors{2}{white}{guidezebra}
\begin{longtable}{|R{0.2754}|R{0.2319}|R{0.4928}|}
\hline
\rowcolor{guidenavy}\hdrcell{You run} & \hdrcell{Prediction} & \hdrcell{Measured on a real Husky + SAM4S} \\ \hline
\endhead
\codesmall{\textquotedbl{}A *= 2;\textquotedbl{}}, \codesmall{MY\_COUNT} = 10, 20, 40 & the loop gets 2x and 4x longer & \textbf{20 cycles (80 samples) per trip} every time; the loop starts near sample 410 and spans \textbf{10.3, 20.4 and 40.3 trips} \\ \hline
\codesmall{\textquotedbl{}A += 2;\textquotedbl{}}, \codesmall{\textquotedbl{}A \textasciicircum{}= 0x55;\textquotedbl{}}, \codesmall{\textquotedbl{}A *= 3;\textquotedbl{}} & maybe different cycle counts & \textbf{all 20 cycles per trip} (20.4, 20.3, 20.3 trips): an add, an xor and a multiply take the same time here, so you \textbf{cannot tell them apart by timing} \\ \hline
\codesmall{\textquotedbl{}A /= 3;\textquotedbl{}}, \codesmall{MY\_\bk{}A\_\bk{}START = \textquotedbl{}0x7FFFFFFF\textquotedbl{}} & the trips shrink, like Step 5 & \textbf{30, 29, 28, 27, 26, 25, 24, 23 cycles} in successive windows \\ \hline
\codesmall{\textquotedbl{}A /= 3;\textquotedbl{}}, \codesmall{MY\_A\_START = \textquotedbl{}0x100\textquotedbl{}} & \codesmall{A} is small from the start & starts at about \textbf{23 cycles} and settles at \textbf{20} after a few trips: once \codesmall{A} reaches 0 the divide is as fast as any other instruction \\ \hline
\end{longtable}}

\vspace{4pt}

\guidefigure{figures/scaling.png}{0.812}{The loop's length scales with \code{MY\_COUNT}: 826, 1628 and 3226 samples for 10, 20 and 40 trips - 80 samples (20 cycles) per trip.}

\guidefigure{figures/divide.png}{0.855}{A constant loop stays at 20 cycles per trip. The divide of a big number starts at 30 and shrinks to 23; the divide of a small number settles at 20 within a few trips.}

\begin{guidecallout}{tip}{GOOD TO KNOW}
\textbf{For the divide loops the trip length is not constant}, so the single-period measurement (\code{measure\_loop})
does not apply. The Step 7 cell notices this and prints the trip length of every window instead (\code{trip length per
window (cycles): [30, 29, 29, .\bk{}.\bk{}.\bk{}]}).
\end{guidecallout}

\guidehone{5  When it does not look like this}

{\small\arrayrulecolor{guidegrid}
\rowcolors{2}{white}{guidezebra}
\begin{longtable}{|R{0.2754}|R{0.3913}|R{0.3333}|}
\hline
\rowcolor{guidenavy}\hdrcell{You see} & \hdrcell{What it means} & \hdrcell{Try} \\ \hline
\endhead
A solid band stuck at +/-0.5, and the printed largest swing is 0.5 & The signal is \textbf{bigger than the ADC can measure} (clipping). It happened on this kit once, after the measurement cable was moved. & Reseat the \textbf{Pos} SMA cable at both ends (finger-tight). If needed \codesmall{scope.gain.db = 15} (or 10) and capture again. \\ \hline
A nearly flat trace, or \codesmall{no repeating loop found} & The measurement cable is not connected, or not on the right connector. & Check the \textbf{Pos} SMA at the Husky and at the CW313 (Lab 0 notebook, Assemble step 4; Linux guide 4.2, steps 2 and 3). \\ \hline
\codesmall{Timeout in OpenADC capture(), no trigger seen} & The scope never saw the trigger. The target is not running the lab firmware (for example, Lab 2\_1B's \codesmall{basic-passwdcheck} is still on it) or \codesmall{JP2} is wrong. & Run the \textbf{Setup cell again} (it reflashes the lab firmware). Check \codesmall{JP2} is on \textbf{T-GPIO4}. \\ \hline
\codesmall{Unable to communicate with found ChipWhisperer} & Another kernel holds the Husky. A failed connection also keeps the USB handle until its kernel stops. & \textbf{Kernel \textgreater{} Shut Down All Kernels}, close other lab tabs, run again. \\ \hline
A number of trips not close to \codesmall{MY\_COUNT} & A window edge, or a weak match. & Re-run once; the green band in the plot must cover the whole loop. \\ \hline
The numbers differ a little from this guide & Normal. Cable seating, gain and the compiler version move amplitudes and, rarely, cycle counts by a cycle or so. & Compare the \textbf{pattern}: constant for \codesmall{*}, \codesmall{+}, \codesmall{\textasciicircum{}}; shrinking for \codesmall{/} with a big \codesmall{A}. \\ \hline
\end{longtable}}

\vspace{4pt}

\guidehone{6  The bundled reference traces}

The captures behind this guide are in
\code{jupyter/\bk{}courses/\bk{}sca101/\bk{}traces/\bk{}hardware\_\bk{}reference/\bk{}} (float32 NumPy files, 5000
samples each). Step 7 of the Explorer plots them before you capture your own. Load one yourself with:

\begin{guidecode}
import numpy as np
x = np.load("traces/hardware_reference/lab2_1A_hw_mul_20.npy")      # 5000 samples, A *= 2; x 20
\end{guidecode}

{\small\arrayrulecolor{guidegrid}
\rowcolors{2}{white}{guidezebra}
\begin{longtable}{|R{0.4493}|R{0.5507}|}
\hline
\rowcolor{guidenavy}\hdrcell{File} & \hdrcell{What it holds} \\ \hline
\endhead
\codesmall{lab2\_\bk{}1A\_\bk{}hw\_\bk{}section0.\bk{}npy} ... \codesmall{section3.npy} & The four programs of the lab (empty, 20 multiplies, loop of 20 multiplies, loop of 20 divides) on a real chip. \\ \hline
\codesmall{lab2\_\bk{}1A\_\bk{}hw\_\bk{}mul\_\bk{}10.\bk{}npy}, \codesmall{mul\_20}, \codesmall{mul\_40} & \codesmall{A *= 2;} repeated 10, 20 and 40 times. \\ \hline
\codesmall{lab2\_\bk{}1A\_\bk{}hw\_\bk{}add\_\bk{}20.\bk{}npy}, \codesmall{xor\_20}, \codesmall{mul3\_20} & \codesmall{A += 2;}, \codesmall{A \textasciicircum{}= 0x55;} and \codesmall{A *= 3;} repeated 20 times. \\ \hline
\codesmall{lab2\_\bk{}1A\_\bk{}hw\_\bk{}div\_\bk{}big\_\bk{}20.\bk{}npy}, \codesmall{div\_small\_20} & \codesmall{A /= 3;} repeated 20 times, with \codesmall{A} starting at \codesmall{0x7FFFFFFF} and at \codesmall{0x100}. \\ \hline
\end{longtable}}

\vspace{4pt}

\textbf{Compare your own capture with a reference:} plot both (\code{plt.\bk{}plot(x); plt.\bk{}plot(mine)}), or run
the notebook's \code{measure\_loop(x)} on each. A healthy capture has the same loop start (about sample 410) and trip
length (80 samples) as the matching reference.

\vspace{10pt}

\guidehthree{Notice}

\guidesmall{This guide accompanies ChipWhisperer course material that references (commercial) training material on
ChipWhisperer.io. You can freely execute and use the lab per the open-source license (including in your own courses if
you distribute similarly), but you must maintain notice about the source: https://www.ChipWhisperer.io.}

\guidesmall{Copyright (C) NewAE Technology Inc., 2015-2020. ChipWhisperer is a trademark of NewAE Technology Inc.,
claimed in all jurisdictions, and registered in at least the United States of America, European Union, and Peoples
Republic of China. Tutorials derived from our open-source work must be released under the associated open-source
license, and notice of the source must be clearly displayed.}

\guidesmall{THE SOFTWARE IS PROVIDED ``AS IS'', WITHOUT WARRANTY OF ANY KIND, EXPRESS OR IMPLIED, INCLUDING BUT NOT
LIMITED TO THE WARRANTIES OF MERCHANTABILITY, FITNESS FOR A PARTICULAR PURPOSE AND NONINFRINGEMENT. IN NO EVENT SHALL
THE AUTHORS OR COPYRIGHT HOLDERS BE LIABLE FOR ANY CLAIM, DAMAGES OR OTHER LIABILITY, WHETHER IN AN ACTION OF
CONTRACT, TORT OR OTHERWISE, ARISING FROM, OUT OF OR IN CONNECTION WITH THE SOFTWARE OR THE USE OR OTHER DEALINGS IN
THE SOFTWARE.}

\end{document}
